\chapter{マシンの全体像}
% full version: chapters/10_synthesis.tex

\pencilfig{10}{\pencilart{10}}{マシンの全体像：データバス、その中のブレーキ、そして上にある4つの制御用放送チャネル。}

ここまで、ニューロンの種類を1つずつ取り上げ、それがマシンに何を加えるかを問うてきた。いよいよ部品を組み合わせるときだ。

できあがったのは3層のマシンである。下には\nt{GLUT}ニューロンの巨大な電話網があり、そこをデータが流れる。その隣に抑制性の\nt{GABA}ニューロンがあり、流れを制御する。何を、いつ、どれだけの大きさで通すかを決めるのだ。そしてそのすべての上に、それぞれ自分のつまみをもつ数個の放送局がある。\nt{DOPA}は、結合のどの変化を残すかを告げる。\nt{SERO}は全体のレベルを保ち、仮説によれば忍耐の視野も決める。\nt{NORA}は探索と決断の間で選ぶ。\nt{ACET}は書き込みと読み出しの間で選ぶ。各放送局は実際には配線1本ではなく小さなチャネルの束だが、それでも860億個のニューロンに対してチャネルはひと握りしかない。

\section*{ひとつの物語}

1つの出来事をマシン全体で追ってみよう。動物はずっと前から知っている。レバーAを押せばジュースがもらえる。ある日、世界が変わる。Aはもう効かず、ジュースをくれるのは、動物がほとんど使ったことのないレバーBになる。

Aのあとにジュースが来ない。ドーパミンは途切れで応え、Aを選んだ結合は弱まる。途切れが繰り返され、誤差がいつもより大きくなる。仮説によれば、ノルアドレナリンが知らせる。世界が変わった。コントラストのつまみが下がり、選択はやわらかくなり、ノイズがBを試すよう後押しすることが増える。アセチルコリンはネットワークを書き込みモードに切り替える。Bを試すと予想外のジュースが出る。ドーパミンが連打し、Bを選んだ結合が強まる。数回試すうちに期待が新しい世界に追いつき、驚きはなくなり、つまみは元の位置に戻る。

注目してほしい。どのつまみも、ほかのつまみが何をしているかを知らない。それぞれが自分の手がかりに応えるだけで、動作の順序はひとりでにできあがる。各ブロックが入力がそろったときに動く非同期回路と同じだ。各つまみが個別にこう振る舞うことは、おおむね測定されている。それらが協調して一緒に働くことは、まだ仮説である。

\section*{ノートパソコンとどこが違うか}

ノートパソコンには共通のクロックがあるが、脳にはない。ノートパソコンの素子は信頼できるが、シナプスは信号の半分以上を失う。ノートパソコンのメモリはプロセッサと別だが、脳では同じ結合の中にある。ノートパソコンは（学習するとしても）誤差逆伝播で学ぶが、脳は局所的な規則と、ブロードキャストされるただ1つの「予想よりよかった」で学ぶ。そして何より、ノートパソコンはプロセッサ1個に数十ワットを使うが、脳はすべてを合わせて20ワットである。

\section*{わかっていないこと}

大脳皮質がカチッの正確なタイミングをどれだけ使っているのか、わかっていない。教師が全員に1人しかいないのに、多層の回路の何十億もの結合をどう調整しているのか、わかっていない。放送局が何を伝えているのかも、完全には理解できていない。共通のクロックなしに、脳の離れた部分がどう仕事を調整しているのかもわからない。そして、同じことを20ワットでこなすマシンを作れる者は、まだ誰もいない。

ここから先、本書はこの中核を越えていく。マシンの入力と出力へ、そのデバッグへ、睡眠へ、そしてチップの上の生きたニューロンで作るマシンへ。

\fullversion{10}
